\chapter{Calf Pain}
\index{Calf Pain}

\section*{Definition and Overview}
Calf pain is a common complaint referring to discomfort in the back of the lower leg, between the knee and the ankle. It can range from a mild muscle ache after exercise to the severe, swollen pain of a deep vein thrombosis, a blood clot in a leg vein.

Because the cause can be harmless or life-threatening, calf pain needs careful assessment, and anyone with sudden, unexplained calf pain or swelling should be checked promptly. Most cases, however, are due to muscle strain or cramp and settle on their own. The key task is to identify the small minority of people whose pain signals a blood clot, a tendon rupture, or a serious injury.

\section*{Epidemiology}
Calf pain is very common and occurs at all ages. Muscle strains and cramps are frequent in people who exercise or who are unaccustomed to activity, and calf cramps at night affect many older adults. Serious causes are less common: deep vein thrombosis occurs in roughly 1 in 1,000 people per year and is more likely after prolonged immobility, surgery, pregnancy, cancer, and in older age. Achilles tendon problems affect athletes and middle-aged people who suddenly increase their activity level.

\section*{Aetiology and Pathophysiology}
Calf pain arises from the muscles, bones, blood vessels, or nerves of the lower leg:
\begin{itemize}
    \item \textbf{Muscle strain or tear:} overstretching of the calf muscles (gastrocnemius or soleus) during sport or a sudden movement
    \item \textbf{Muscle cramp:} a sudden, involuntary contraction, often at night or after exercise, linked to dehydration, fatigue, or nerve irritation
    \item \textbf{Deep vein thrombosis (DVT):} a blood clot in a deep leg vein, usually after prolonged immobility, surgery, or in people with an inherited clotting tendency
    \item \textbf{Achilles tendon problems:} tendinopathy or a torn Achilles tendon from overuse or sudden loading
    \item \textbf{Compartment syndrome:} a dangerous build-up of pressure within the leg after a severe injury
    \item \textbf{Nerve causes:} sciatica or a trapped nerve in the back sending pain into the calf
    \item \textbf{Other causes:} a Baker's cyst (a fluid-filled swelling behind the knee), peripheral artery disease, or a fracture
\end{itemize}

\section*{Clinical Features}
The features vary with the cause, and the pattern is important:
\begin{itemize}
    \item \textbf{Muscle strain:} a sudden, sharp pain during activity, followed by tenderness, bruising, and pain on stretching or rising onto the toes
    \item \textbf{Cramp:} a sudden, intense, knotted pain that relaxes over minutes
    \item \textbf{DVT:} a swollen, warm, tender calf, often with a deep ache and sometimes redness; about half of people with a clot have no symptoms
    \item \textbf{Achilles rupture:} a sudden snap or kick in the calf with severe pain and difficulty pushing off the foot
    \item \textbf{Sciatica:} burning or shooting pain travelling from the lower back or buttock down the back of the leg
    \item \textbf{Arterial disease:} calf pain or cramping brought on by walking and relieved by rest, known as claudication
\end{itemize}

\section*{Differential Diagnosis}
Because many conditions cause calf pain, the doctor distinguishes between them:
\begin{itemize}
    \item Deep vein thrombosis versus muscle strain or a ruptured Baker's cyst
    \item Achilles tendon tear versus a calf muscle tear
    \item Sciatica or a trapped nerve versus simple muscle pain
    \item Peripheral arterial disease versus a nerve or muscle cause of cramp
    \item Compartment syndrome versus a simple bruise or strain after injury
    \item A bone fracture versus a soft-tissue injury
\end{itemize}

\section*{When to Seek Medical Care}
People should see a doctor promptly for any new, unexplained calf pain or swelling, especially if they have risk factors for a blood clot, such as recent surgery, immobility, pregnancy, or cancer.

\textbf{Contact your local emergency services for:}
\begin{itemize}
    \item Sudden chest pain or difficulty breathing, which may indicate a pulmonary embolism from a clot
    \item A leg that is pale, cold, numb, or unable to move
    \item Severe pain after an injury, or a limb that looks deformed
    \item A hot, swollen, extremely tense calf after injury, which may indicate compartment syndrome
\end{itemize}

\section*{Diagnosis and Investigations}
Most calf pain is diagnosed from the history and examination, with tests used when a serious cause is possible:
\begin{itemize}
    \item \textbf{History:} how the pain started, activity before it, risk factors for clots, and any associated symptoms
    \item \textbf{Examination:} checking for tenderness, swelling, warmth, bruising, and comparing both legs
    \item \textbf{Ultrasound:} the main test for deep vein thrombosis, looking directly for a clot in the leg veins
    \item \textbf{D-dimer blood test:} used to help rule out a clot when the pre-test risk is low
    \item \textbf{X-ray or other imaging:} if a fracture or bone problem is suspected
    \item \textbf{Ultrasound or MRI:} may be used to assess the Achilles tendon or a suspected Baker's cyst
\end{itemize}

\section*{Management}
Management depends on the cause and ranges from self-care to urgent hospital treatment:
\begin{itemize}
    \item \textbf{Muscle strain:} rest, ice, gentle compression, elevation, and simple pain relief such as paracetamol, followed by gentle stretching once the acute pain settles
    \item \textbf{Cramp:} gentle stretching of the calf, warmth, and attention to hydration and footwear
    \item \textbf{DVT:} anticoagulant (blood-thinning) medicines, usually continued for several months, with early walking while on treatment
    \item \textbf{Achilles tear:} a walking boot, physiotherapy, and sometimes surgery
    \item \textbf{Baker's cyst:} rest, anti-inflammatory treatment, and treatment of the underlying knee problem
    \item \textbf{Sciatica:} simple pain relief, activity as tolerated, and physiotherapy, with specialist care if severe
    \item \textbf{Physiotherapy:} a graded return to activity and exercises to restore strength and prevent recurrence
\end{itemize}

\section*{Complications}
\begin{itemize}
    \item Untreated DVT can break off and travel to the lungs, causing a life-threatening pulmonary embolism
    \item Long-term swelling and skin changes in the leg after a DVT, known as post-thrombotic syndrome
    \item Recurrent strains and weakness if the calf is not rehabilitated properly
    \item Chronic tendinopathy or re-tear of the Achilles tendon
    \item Chronic pain or nerve damage from a missed fracture or compartment syndrome
\end{itemize}

\section*{Prognosis and Outcomes}
Most calf pain from muscle strain settles within a few days to weeks, with gradual return to activity. Cramps are usually harmless and respond to simple measures. A DVT is very treatable with blood-thinning medicines, and serious complications are uncommon with early treatment. Achilles tendon tears take longer, often several months, to regain strength. Overall, the outlook is good provided that serious causes are identified and treated early.

\section*{Prevention and Self-Care}
\begin{itemize}
    \item Warm up and stretch before exercise, and build up activity gradually
    \item Stay well hydrated and replace salt lost through heavy sweating
    \item Wear supportive, well-fitted footwear
    \item Stretch the calf gently before bed if you get night cramps, and keep the leg warm
    \item Avoid long periods of sitting or lying still; on long flights or car trips, stand, walk, and move your legs regularly
    \item After surgery or a hospital stay, follow advice about moving early to reduce the risk of blood clots
    \item See a doctor promptly for any unexplained, persistent, or worsening calf pain
\end{itemize}

\section*{Sources and Further Reading}
The following sources provide reliable, accessible information on calf pain and its causes:
\begin{itemize}
    \item NHS. \textit{Leg pain.} https://www.nhs.uk/conditions/leg-pain/
    \item Mayo Clinic. \textit{Calf pain.} https://www.mayoclinic.org/
    \item MedlinePlus. \textit{Calf pain.} https://medlineplus.gov/
    \item MSD Manual. \textit{Deep vein thrombosis.} https://www.msdmanuals.com/
    \item NICE. \textit{Venous thromboembolic diseases.} https://www.nice.org.uk/
\end{itemize}
